\section{Hardware implementation}
\label{sec:hardware}

\subsection{Premium-qubit selection}
\label{sec:premium}

To run on hardware we screen candidate chains (rings) with two
quantities: the stabilizer mean, probing two-qubit gate quality, and
the readout fidelity, probing measurement quality.
The ideal output is the cluster state
$|C\rangle=\prod_{\langle ij\rangle}CZ_{ij}|+\rangle^{\otimes n}$,
where $\langle ij\rangle$ runs over the edges of the candidate subgraph;
its stabilizers $S_i$ (two-body $X_0Z_1$, $Z_{n-2}X_{n-1}$ at open ends,
three-body $Z_{i-1}X_iZ_{i+1}$ otherwise) have expectation $+1$ in the
absence of errors, so deviations of $\langle S_i\rangle$ from $1$
directly reflect $CZ$ gate errors.
Each candidate is characterized by $4$ benchmark circuits: \texttt{stab\_g0}
and \texttt{stab\_g1} measure the even- and odd-index stabilizers separately
(the two groups anticommute and cannot be read out together), while
\texttt{allzero} and \texttt{allone} contain no $CZ$ gates and yield the readout
fidelity $F_{\rm ro}$ from the mean of $P(00\cdots0)$ and
$P(11\cdots11)$. The score
\begin{equation}
{\rm score}=0.8\,\bar{S}+0.2\,F_{\rm ro},
\end{equation}
where $\bar{S}=n^{-1}\sum_i\langle S_i\rangle$ is the stabilizer mean,
ranks the candidates, weighting two-qubit gate quality over readout.
Cluster states~\cite{raussendorf2001oneway}
are the ideal outputs whose stabilizers we measure; measurement-based
computation on such states is reviewed in
Ref.~\cite{briegel2009measurement}.
The $8$-chain and $10$-ring benchmarks differ (the ring carries one
extra closing $CZ$ and all-three-body stabilizers), so rankings are only
compared within each geometry, never across. In contrast to generic
randomized benchmarks, this screening probes exactly the stabilizers of
the states the hardware will prepare.
Hardware characterization around this screening takes several
complementary forms: Pauli-frame randomization has been demonstrated on
a superconducting qubit and validated by gate-set
tomography~\cite{ware2021experimental}, while the analysis of
randomized benchmarking under time-correlated dephasing shows where the
standard Markovian decay model stops
applying~\cite{qi2021randomized}. Shadow-based Hamiltonian measurement
offers a measurement-frugal
alternative~\cite{hadfield2022measurements,huang2021derandomization}, and
shadow tomography extends to quantum
channels~\cite{kunjummen2023shadow}.

The resulting rankings and selected qubits are reported in the Results and
Discussion feature. This Methods section defines only the screening
protocol and its within-geometry comparison rule.

\subsection{Premium-qubit rankings}
\label{sec:rankings}

The screening reports three rankings on the Baihua processor: (a) top-ten
$8$-chain candidates, (b) top-ten $10$-ring candidates, and (c) top-ten
$8$-subchains within the champion ring. The specific qubit tables follow
the latest S04 manifest; until the D07 figures land the slots below are
reserved and no ranking is asserted.

\begin{figure}[htbp]
\centering
\fbox{\parbox{0.9\columnwidth}{\centering TODO(GAP-001): D07a ranking figure pending.}}\\(a)
\fbox{\parbox{0.9\columnwidth}{\centering TODO(GAP-001): D07b ranking figure pending.}}\\(b)
\fbox{\parbox{0.9\columnwidth}{\centering TODO(GAP-001): D07c ranking figure pending.}}\\(c)
\caption{Premium-qubit ranking placeholders. \label{fig:D07}}
\end{figure}
